\documentclass{amsart}
\usepackage{amsmath,amssymb,amsthm,mathtools,hyperref}
\newtheorem{theorem}{Theorem}
\newtheorem{lemma}{Lemma}
\newtheorem{proposition}{Proposition}
\theoremstyle{definition}
\newtheorem{definition}{Definition}
\newtheorem{remark}{Remark}

\begin{document}

\title{Representability of uniform matroids}
\maketitle

Throughout, $a,b$ are integers with $0\le a\le b$, and $F$ is a field. We write
$|F|$ for the cardinality of $F$ (possibly infinite). Recall that $U_{a,b}$ is the
matroid on $E=\{1,\dots,b\}$ whose independent sets are exactly the subsets of $E$
of size at most $a$; in particular its rank is $a$ and the bases are exactly the
$a$-element subsets.

\section{Statement of the answer}

\begin{theorem}\label{thm:main}
Let $0\le a\le b$.
\begin{enumerate}
\item\textbf{(Trivial range)} If $a\in\{0,1,b-1,b\}$, then $U_{a,b}$ is
representable over \emph{every} field $F$.
\item\textbf{(Main range)} Suppose $2\le a\le b-2$. Then:
\begin{enumerate}
\item If $|F|\ge b-1$ (in particular for every infinite field), then $U_{a,b}$ is
representable over $F$.
\item If $|F|\le b-3$, then $U_{a,b}$ is \emph{not} representable over $F$.
\item In the boundary case $|F|=b-2$ (a finite field $F=\mathbb F_q$ with
$b=q+2$): $U_{a,q+2}$ is representable over $\mathbb F_q$ if and only if $q$ is
even and $a\in\{3,\,q-1\}$.
\end{enumerate}
\end{enumerate}
\end{theorem}

In summary, writing $q=|F|$ when $F$ is finite (and $q=\infty$ when $F$ is
infinite, with the convention $q+1=q+2=\infty$):
\[
U_{a,b}\text{ is representable over }F\iff
\begin{cases}
\text{always}, & a\in\{0,1,b-1,b\},\\[2pt]
b\le q+1, & 2\le a\le b-2,\ q\ \text{odd or}\ a\notin\{3,q-1\},\\[2pt]
b\le q+2, & 2\le a\le b-2,\ q\ \text{even and}\ a\in\{3,q-1\}.
\end{cases}
\]

The ``only if'' direction of part (2)(c), together with the sharp universal bound
$b\le q+2$, is the maximum-length theorem for MDS codes (the MDS conjecture),
proved by Ball for fields of prime order and by Ball--De~Beule in general. We
isolate exactly where this external input is used (Theorem~\ref{thm:mds}); every
other statement, including all the ``representable'' directions, is proved here
from scratch.

\section{Reformulation in terms of arcs and MDS codes}

\begin{lemma}\label{lem:reform}
Let $1\le a\le b$. The matroid $U_{a,b}$ is representable over $F$ if and only if
there exist vectors $v_1,\dots,v_b\in F^{a}$ such that every $a$ of them are
linearly independent (i.e.\ form a basis of $F^a$).
\end{lemma}

\begin{proof}
($\Leftarrow$) Suppose $v_1,\dots,v_b\in F^a$ are such that every $a$-subset is a
basis. We claim these vectors represent $U_{a,b}$, i.e.\ for $I\subseteq E$ the set
$\{v_i:i\in I\}$ is linearly independent over $F$ iff $|I|\le a$.

If $|I|\le a$, extend $I$ to a set $J\supseteq I$ with $|J|=a$ (possible since
$a\le b$). By hypothesis $\{v_j:j\in J\}$ is a basis of $F^a$, hence linearly
independent, so its subfamily $\{v_i:i\in I\}$ is linearly independent. If
$|I|>a$, then $\{v_i:i\in I\}$ consists of more than $a=\dim F^a$ vectors and is
therefore linearly dependent. Thus the family represents exactly $U_{a,b}$.

($\Rightarrow$) Suppose $U_{a,b}$ is represented by vectors $w_1,\dots,w_b$ in some
$F$-vector space $W$. Since the matroid has rank $a$, the subspace
$V=\mathrm{span}_F(w_1,\dots,w_b)$ has dimension $a$: any basis of $U_{a,b}$ (an
$a$-subset) is linearly independent and is a maximal independent set, hence spans
$V$, so $\dim V=a$. Choose an $F$-linear isomorphism $\phi:V\to F^a$ and set
$v_i=\phi(w_i)\in F^a$. Linear (in)dependence is preserved by an isomorphism, so
the $v_i$ also represent $U_{a,b}$; in particular every $a$-subset of the $v_i$ is
a basis of $U_{a,b}$, hence linearly independent, hence a basis of $F^a$.
\end{proof}

\begin{remark}\label{rem:projective}
For $2\le a\le b-2$, in any representation as in Lemma~\ref{lem:reform} the
vectors $v_i$ are nonzero (a single $v_i$ is independent because $1\le a$) and
pairwise non-proportional (two proportional vectors are dependent, but any
$2\le a$ of the $v_i$ are independent). Hence they determine $b$ distinct points
$[v_1],\dots,[v_b]$ in the projective space $\mathrm{PG}(a-1,F)=\mathbb P(F^a)$,
no $a$ of which lie on a common hyperplane through the origin; such a configuration
is a \emph{$b$-arc in $\mathrm{PG}(a-1,F)$}. Conversely a $b$-arc lifts to such
vectors. Equivalently, the $a\times b$ matrix with columns $v_i$ is a generator
matrix of a linear $[b,a]$ MDS code over $F$. Thus, for $2\le a\le b-2$,
\[
U_{a,b}\ \text{representable over }F
\iff \text{a }b\text{-arc exists in }\mathrm{PG}(a-1,F)
\iff \text{an }[b,a]\text{ MDS code over }F\text{ exists.}
\]
\end{remark}

\section{The trivial range $a\in\{0,1,b-1,b\}$}

\begin{proposition}\label{prop:trivial}
If $a\in\{0,1,b-1,b\}$ then $U_{a,b}$ is representable over every field $F$.
\end{proposition}

\begin{proof}
We treat the four (possibly overlapping) cases.

\emph{Case $a=0$.} Then $U_{0,b}$ has independent sets exactly $\{\emptyset\}$.
Take $v_1=\dots=v_b=0\in F^1$. A subset $\{v_i:i\in I\}$ is linearly independent
over $F$ iff $I=\emptyset$ (any nonempty family containing $0$ is dependent). This
is exactly $U_{0,b}$.

\emph{Case $a=b$.} Then $U_{b,b}$ is the free matroid. Take $v_1,\dots,v_b$ the
standard basis of $F^b$; every subset of a basis is independent, matching
$U_{b,b}$.

\emph{Case $a=1$ (with $b\ge1$).} $U_{1,b}$ has independent sets the subsets of
size $\le1$. Take $v_1=\dots=v_b=1\in F^1$. A singleton $\{v_i\}$ is independent
($1\ne0$); any set of $\ge2$ of these is dependent (e.g.\
$1\cdot v_i+(-1)\cdot v_j=0$, a nontrivial relation for $i\ne j$); and
$\emptyset$ is independent. So the family represents $U_{1,b}$, over every field
including $\mathbb F_2$.

\emph{Case $a=b-1$ (with $b\ge1$).} In $F^{a}=F^{b-1}$ take $v_i=e_i$ for
$1\le i\le b-1$ (standard basis) and $v_b=e_1+\dots+e_{b-1}$ (all-ones vector). By
Lemma~\ref{lem:reform} it suffices to show that each of the $b$ subsets of size
$a=b-1$ is a basis of $F^{b-1}$. The subsets are $S_k=\{v_1,\dots,v_b\}\setminus
\{v_k\}$, $k=1,\dots,b$.
\begin{itemize}
\item $S_b=\{e_1,\dots,e_{b-1}\}$ is the standard basis.
\item For $1\le k\le b-1$: the matrix $M_k$ with columns $e_1,\dots,e_{k-1},
e_{k+1},\dots,e_{b-1},\mathbf 1$ (where $\mathbf 1=v_b$) has determinant equal to
the $k$-th coordinate of $\mathbf 1$, up to sign. Indeed, expand $\det M_k$ along
the last column $\mathbf 1$: the cofactor of the $i$-th entry of $\mathbf 1$ is
$\pm$ the determinant of the matrix obtained by deleting row $i$ and the last
column, whose columns are $b-2$ distinct standard basis vectors from
$\{e_1,\dots,e_{b-1}\}\setminus\{e_k\}$; this minor is nonzero only when the
missing index of the row equals the missing index $k$ of the columns, i.e.\ only
for $i=k$, and then the minor equals $\pm1$. Hence $\det M_k=\pm(\mathbf 1)_k=\pm
1\ne0$, so $S_k$ is a basis.
\end{itemize}
Thus every $(b-1)$-subset is a basis, and $U_{b-1,b}$ is representable over every
$F$.
\end{proof}

\section{Sufficiency in the main range: $|F|\ge b-1$}

\begin{proposition}\label{prop:suff}
Let $2\le a\le b-2$ and suppose $|F|\ge b-1$. Then $U_{a,b}$ is representable over
$F$.
\end{proposition}

\begin{proof}
Since $|F|\ge b-1$, choose $b-1$ \emph{distinct} elements
$t_1,\dots,t_{b-1}\in F$. Define $b$ vectors in $F^a$:
\[
v_i=\bigl(1,\;t_i,\;t_i^2,\;\dots,\;t_i^{\,a-1}\bigr)^{\!\top}
\quad(1\le i\le b-1),
\qquad
v_b=(0,0,\dots,0,1)^{\!\top}\in F^a .
\]
By Lemma~\ref{lem:reform} it suffices to show every $a$-subset of
$\{v_1,\dots,v_b\}$ has nonzero determinant.

\emph{Subsets not containing $v_b$.} For $1\le i_1<\dots<i_a\le b-1$ the
$a\times a$ matrix with columns $(1,t_{i_j},\dots,t_{i_j}^{a-1})^\top$ is a
Vandermonde matrix, of determinant
\[
\prod_{1\le r<s\le a}\bigl(t_{i_s}-t_{i_r}\bigr)\ne0,
\]
since the $t_{i_j}$ are pairwise distinct.

\emph{Subsets containing $v_b$.} For $1\le i_1<\dots<i_{a-1}\le b-1$, let $M$ have
columns $v_{i_1},\dots,v_{i_{a-1}},v_b$. The last column is $e_a$, so expanding
along it gives $\det M=\det M'$, where $M'$ is the $(a-1)\times(a-1)$ matrix
obtained by deleting the last row and last column. Then $M'$ has columns
$(1,t_{i_j},\dots,t_{i_j}^{\,a-2})^\top$, a Vandermonde matrix in
$t_{i_1},\dots,t_{i_{a-1}}$, so
\[
\det M=\det M'=\prod_{1\le r<s\le a-1}\bigl(t_{i_s}-t_{i_r}\bigr)\ne0
\]
(the empty product equals $1$ when $a=2$). Hence every $a$-subset is a basis, and
$U_{a,b}$ is representable over $F$.
\end{proof}

\begin{remark}
Geometrically, $[v_1],\dots,[v_{b-1}]$ lie on the rational normal curve
$\{[1:t:\dots:t^{a-1}]\}$ of $\mathrm{PG}(a-1,F)$, and $[v_b]$ is its point at
infinity; together they form a $b$-arc whenever $|F|\ge b-1$.
\end{remark}

\section{Duality}

\begin{lemma}\label{lem:dual}
For $0\le a\le b$, the dual matroid of $U_{a,b}$ is $U_{b-a,b}$, and $U_{a,b}$ is
representable over a field $F$ if and only if $U_{b-a,b}$ is.
\end{lemma}

\begin{proof}
The bases of $U_{a,b}$ are the $a$-subsets of $E$; the bases of the dual
$U_{a,b}^{*}$ are by definition their complements, the $(b-a)$-subsets of $E$.
The matroid whose bases are exactly the $(b-a)$-subsets is $U_{b-a,b}$. Hence
$U_{a,b}^{*}=U_{b-a,b}$.

For representability, assume $U_{a,b}$ is representable over $F$; if $a\in\{0,b\}$
the claim follows from Proposition~\ref{prop:trivial} applied to $U_{b-a,b}$, so
assume $1\le a\le b-1$. By Lemma~\ref{lem:reform} take a full-rank $a\times b$
matrix $G=[v_1\,|\cdots|\,v_b]$ over $F$ representing $U_{a,b}$, i.e.\ every
$a$-subset of columns is a basis. After a change of coordinates of $F^a$
(an invertible row operation, which preserves all column (in)dependencies) we may
assume the first $a$ columns form the identity, so $G=[\,I_a\,|\,A\,]$ for some
$a\times(b-a)$ matrix $A$. Set $H=[\,-A^{\top}\,|\,I_{b-a}\,]$, a full-rank
$(b-a)\times b$ matrix with $GH^{\top}=-A+A=0$, whose rows therefore span
$\ker G$. We claim $H$ represents $U_{b-a,b}$, i.e.\ every $(b-a)$-subset of
columns of $H$ is a basis of $F^{b-a}$.

Let $S\subseteq E$ with $|S|=b-a$ and $T=E\setminus S$, $|T|=a$. We show
$H_S$ (the columns of $H$ indexed by $S$) is invertible iff $G_T$ is. By the
generalized Laplace/complementary-minor identity for the pair $(G,H)$ with
$G=[I_a\mid A]$, $H=[-A^\top\mid I_{b-a}]$, one has the exact relation
\[
\det H_S=\pm\det G_T
\]
for every such complementary pair $(S,T)$. Concretely, this is the standard fact
that for a matrix in the form $[\,I_a\mid A\,]$ the $a\times a$ minor on a column
set $T$ equals (up to sign) the complementary $(b-a)\times(b-a)$ minor
$\det H_S$ of $[-A^\top\mid I_{b-a}]$; both reduce, by Laplace expansion, to the
same minor of $A$ (the submatrix of $A$ on the rows indexed by $T\cap\{1,\dots,a\}$
and columns indexed by $S\cap\{a+1,\dots,b\}$), once one accounts for the identity
blocks. Hence $\det H_S\ne0\iff\det G_T\ne0$. Therefore every $(b-a)$-subset of
columns of $H$ is a basis iff every $a$-subset of columns of $G$ is a basis; the
latter holds, so $H$ represents $U_{b-a,b}$. The converse direction is identical
with the roles of $G,H$ exchanged (since $(U_{a,b}^{*})^{*}=U_{a,b}$).
\end{proof}

\section{Necessity in the main range and the sharp MDS input}

\begin{lemma}[Elementary rank-$2$ bound]\label{lem:rank2}
Let $F=\mathbb F_q$ be finite. If $U_{2,m}$ is representable over $\mathbb F_q$,
then $m\le q+1$. Moreover $U_{2,m}$ is representable over $\mathbb F_q$ if and
only if $m\le q+1$.
\end{lemma}

\begin{proof}
By Lemma~\ref{lem:reform}, representability of $U_{2,m}$ means there are vectors
$u_1,\dots,u_m\in\mathbb F_q^{2}$, every two independent, i.e.\ nonzero and
pairwise non-proportional. The number of one-dimensional subspaces of
$\mathbb F_q^2$ is $\tfrac{q^2-1}{q-1}=q+1$, and the $u_i$ lie in pairwise
distinct ones; hence $m\le q+1$. Conversely, if $m\le q+1$ then either $m\le2$
(trivial range, Proposition~\ref{prop:trivial}) or $3\le m\le q+1$, and
Proposition~\ref{prop:suff} with $a=2,b=m$ (note $|F|=q\ge m-1=b-1$) gives a
representation.
\end{proof}

\begin{theorem}[Maximum length of MDS codes; arc bound]\label{thm:mds}
Let $\mathbb F_q$ be a finite field and $2\le a\le b-2$.
\begin{enumerate}
\item \emph{(Universal bound.)} Every $b$-arc in $\mathrm{PG}(a-1,q)$ satisfies
$b\le q+2$.
\item \emph{(Refined bound.)} Moreover $b\le q+1$ unless $q$ is even and
$a\in\{3,\,q-1\}$, in which case the maximum arc length is exactly $q+2$.
\end{enumerate}
\end{theorem}

This is the resolution of the MDS conjecture: the universal bound (1) and the
refinement (2) for $q$ of prime order are theorems of Ball, and the full
refinement for prime powers is due to Ball and Ball--De~Beule. We use it as a
cited input. By Remark~\ref{rem:projective} it translates into the following
statement about uniform matroids, which is the necessity half of our theorem.

\begin{proposition}\label{prop:nec}
Let $2\le a\le b-2$ and $F=\mathbb F_q$.
\begin{enumerate}
\item If $U_{a,b}$ is representable over $\mathbb F_q$, then $b\le q+2$.
\item If $b=q+2$, then $U_{a,b}$ is representable over $\mathbb F_q$ if and only
if $q$ is even and $a\in\{3,q-1\}$.
\end{enumerate}
\end{proposition}

\begin{proof}
(1) If $U_{a,b}$ is representable then by Remark~\ref{rem:projective} there is a
$b$-arc in $\mathrm{PG}(a-1,q)$, so $b\le q+2$ by Theorem~\ref{thm:mds}(1).

(2) ($\Rightarrow$) A representation of $U_{a,q+2}$ gives a $(q+2)$-arc in
$\mathrm{PG}(a-1,q)$ (Remark~\ref{rem:projective}); by Theorem~\ref{thm:mds}(2)
this forces $q$ even and $a\in\{3,q-1\}$.

($\Leftarrow$) Suppose $q$ is even and $a=3$, $b=q+2$. Take the $q+2$ vectors of
$\mathbb F_q^3$
\[
v_t=(1,\,t,\,t^2)^{\top}\ (t\in\mathbb F_q),\qquad
v_\infty=(0,0,1)^{\top},\qquad
v_\star=(0,1,0)^{\top}.
\]
(The points $[v_t]$ and $[v_\infty]$ form the conic
$x_0x_2=x_1^2$ together with its point at infinity, and $[v_\star]$ is its
nucleus.) By Lemma~\ref{lem:reform} we check that every $3$-subset is a basis of
$\mathbb F_q^3$, by computing determinants of the $3\times3$ matrices whose columns
are the chosen vectors. There are four shapes of $3$-subset:
\begin{itemize}
\item three conic points $v_{t_1},v_{t_2},v_{t_3}$ ($t_i$ distinct):
determinant $=\prod_{r<s}(t_s-t_r)\ne0$ (Vandermonde);
\item two conic points and $v_\infty$: $\det\!\begin{pmatrix}1&1&0\\
t_1&t_2&0\\ t_1^2&t_2^2&1\end{pmatrix}=t_2-t_1\ne0$;
\item two conic points and $v_\star$: $\det\!\begin{pmatrix}1&1&0\\
t_1&t_2&1\\ t_1^2&t_2^2&0\end{pmatrix}=-(t_2^2-t_1^2)=(t_2-t_1)^2\ne0$ (in
characteristic $2$);
\item one conic point and $v_\infty,v_\star$: $\det\!\begin{pmatrix}1&0&0\\
t_1&0&1\\ t_1^2&1&0\end{pmatrix}=-1\ne0$.
\end{itemize}
(Each determinant is nonzero precisely because the $t_i$ are distinct.) Hence
every $3$-subset is a basis, so $U_{3,q+2}$ is representable over $\mathbb F_q$
for $q$ even. Finally, for $q$ even and $a=q-1$, note $b-a=(q+2)-(q-1)=3$; by
duality (Lemma~\ref{lem:dual}) $U_{q-1,q+2}$ is representable over $\mathbb F_q$
iff $U_{3,q+2}$ is, which we have just shown. This completes both directions.
\end{proof}

\section{Proof of the main theorem}

\begin{proof}[Proof of Theorem~\ref{thm:main}]
Part (1) is Proposition~\ref{prop:trivial}. Assume now $2\le a\le b-2$.

(2)(a) If $|F|\ge b-1$, Proposition~\ref{prop:suff} gives a representation (for
infinite $F$, $|F|\ge b-1$ holds trivially).

(2)(b) Suppose $|F|\le b-3$ and $U_{a,b}$ is representable over $F$. Then $F$ is
finite, $F=\mathbb F_q$ with $q\le b-3$, i.e.\ $b\ge q+3$. By
Proposition~\ref{prop:nec}(1), representability forces $b\le q+2$, a
contradiction. Hence $U_{a,b}$ is not representable over $F$.

(2)(c) If $|F|=b-2$ then $F=\mathbb F_q$ with $b=q+2$, and
Proposition~\ref{prop:nec}(2) gives: representable iff $q$ even and
$a\in\{3,q-1\}$.

For finite $F=\mathbb F_q$, the three cases $|F|\ge b-1\ (\!\iff b\le q+1)$,
$|F|=b-2\ (\!\iff b=q+2)$, and $|F|\le b-3\ (\!\iff b\ge q+3)$ are exhaustive and
mutually exclusive, and over any infinite field $|F|\ge b-1$ always holds. This
yields exactly the displayed characterization and completes the proof.
\end{proof}

\section{The role of the external input, and unconditional cases}

The only non-self-contained ingredient is Theorem~\ref{thm:mds}, used solely to
obtain the sharp bound $b\le q+2$ (Proposition~\ref{prop:nec}(1)) and to certify
that $b=q+2$ occurs \emph{only} in the two even-characteristic cases
(the ``$\Rightarrow$'' of Proposition~\ref{prop:nec}(2)). All ``representable''
directions are constructive and proved here in full
(Propositions~\ref{prop:trivial},~\ref{prop:suff},~\ref{prop:nec}($\Leftarrow$);
Lemmas~\ref{lem:rank2},~\ref{lem:dual}).

Consequently the answer is unconditional in the following ubiquitous cases:
\begin{itemize}
\item $F$ infinite, or $a\in\{0,1,b-1,b\}$: representable.
\item $a=2$ (or, by duality, $a=b-2$): the elementary
Lemma~\ref{lem:rank2} shows $U_{2,b}$ is representable over $\mathbb F_q$ iff
$b\le q+1$, with no MDS input needed; the even-$q$ exception of
Theorem~\ref{thm:mds} requires $a\in\{3,q-1\}$, which excludes $a=2$ unless
$q-1=2$, i.e.\ $q=3$ (odd, hence no exception). So for $a=2$ the threshold is
exactly $b\le q+1$, unconditionally.
\item $F=\mathbb F_q$ with $q$ \emph{prime}: Theorem~\ref{thm:mds} is then Ball's
theorem for prime fields, and the even-$q$ exception is vacuous except $q=2$,
where $a\in\{3,q-1\}=\{3,1\}$ forces $a=3,b=4$, lying in the trivial range
$a=b-1$ (and indeed $U_{3,4}$ is representable over $\mathbb F_2$). Hence for $q$
prime: $U_{a,b}$ ($2\le a\le b-2$) is representable over $\mathbb F_q$ iff
$b\le q+1$.
\end{itemize}

\end{document}
